\documentclass{article}
\usepackage{amsmath}
\usepackage{amssymb}

\begin{document}

\[ \text{Exercice N° 8:} \]
\[ \alpha \begin{pmatrix} 1 \\ 1 \\ 1 \end{pmatrix} + \beta \begin{pmatrix} 2 \\ 2 \\ 2 \end{pmatrix} = \begin{pmatrix} a \\ 3 \\ 7 \end{pmatrix} \]
\[ \begin{cases} \alpha + 2\beta = a \\ \alpha + 2\beta = 3 \\ \alpha + 2\beta = 7 \end{cases} \]
\[ \alpha \begin{pmatrix} 1 \\ 1 \\ 1 \end{pmatrix} + \beta \begin{pmatrix} 2 \\ 2 \\ 2 \end{pmatrix} + \gamma \begin{pmatrix} a \\ 3 \\ -7 \end{pmatrix} = \vec{0} \]
\[ \begin{cases} \alpha + 2\beta + a\gamma = 0 \\ \alpha + 2\beta + 3\gamma = 0 \\ \alpha + 2\beta - 7\gamma = 0 \end{cases} \]
\[ \begin{cases} \alpha + 2\beta + a\gamma = 0 \\ \gamma(a-3) = 0 \\ \gamma(a-7) = 0 \end{cases} \]
\[ 10\gamma = 0 \implies \gamma = 0 \]
\[ a = \mathbb{R} \text{ car } \gamma = 0 \]
\[ \text{toujours nul} \]

\end{document}